\subsection{半双线性型的张量积与外幂。}

本节中，A 表示一个交换环。因此，两个 A-模的积上的双线性型是半双线性型的一个特例。我们用 J 表示 A 的一个自同构，用 J′ 表示其逆。

设 \(\mathbf{E}_i\) （ \(i = 1, \ldots, m\) ）是一些 A-模。映射

\[
(x _ {1}, \dots , x _ {m}) \rightarrow x _ {1} \otimes \dots \otimes x _ {m}
\]

（从 \(\prod_{i=1}^{m}\mathrm{E}_{i}^{j}\) 到 \((\bigotimes_{i=1}^{m}\mathrm{E}_{i})^{j}\) 中，\(x_{i}\in\mathrm{E}_{i}^{j}\) ）（参看第 2 小节，定义 5）显然是 A-重线性的；于是它定义了（第 III 章，第 1 节，第 7 小节）一个 A-线性映射 \(f\) （从 \(\bigotimes_{i}\mathrm{E}_{i}^{\mathrm{J}}\) 到 \((\bigotimes_{i}\mathrm{E}_{i})^{\mathrm{J}}\) 中）；这个映射把 \(x_{1}\otimes\cdots\otimes x_{m}\) （其中记号 \(\otimes\) 表示 \(\bigotimes_{i}\mathrm{E}_{i}^{\mathrm{J}}\) 中的张量积）映为 \(x_{1}\otimes\cdots\otimes x_{m}\) （其中记号 \(\otimes\) 表示 \((\bigotimes_{i}\mathrm{E}_{i})^{\mathrm{J}}\) 中的张量积）。于是 \(f\) 是 \(\bigotimes_{i}\mathrm{E}_{i}^{\mathrm{J}}\) 到 \((\bigotimes_{i}\mathrm{E}_{i})^{\mathrm{J}}\) 上的一个同构。我们将借助这个同构把这两个模等同起来。

类似地，设 E 是一个 A-模。映射

\[
(x _ {1}, \dots , x _ {m}) \rightarrow x _ {1} \wedge \dots \wedge x _ {m}
\]

（从 \((\mathrm{E}^{\mathfrak{J}})^{m}\) 到 \((\wedge^{\mathfrak{m}}\mathrm{E})^{\mathfrak{J}}\) 中）显然是 A-重线性的且交错的。于是它定义了一个 A-线性映射 \(f\) （从 \(\wedge^{\mathfrak{m}}\mathrm{E}^{\mathfrak{J}}\) 到 \((\wedge^{\mathfrak{m}}\mathrm{E})^{\mathfrak{J}}\) 中），后者显然是一个同构。我们将借助这个同构把 \(\wedge^{\mathfrak{m}}\mathrm{E}^{\mathfrak{J}}\) 与 \((\wedge^{\mathfrak{m}}\mathrm{E})^{\mathfrak{J}}\) 等同起来。

设 \(x'\) 是 \(\mathbf{E}^*\) （\(\mathbf{E}\) 的对偶）的一个元素。映射 \(x \to \langle x, x' \rangle^{\mathfrak{J}}\) （ \(x \in \mathbf{E}\) ）是一个元素 \(x'^{\mathfrak{J}}\) （属于 \((\mathrm{E}^{\mathfrak{J}})^*\) ），且立即看出 \(x' \to x'^{\mathfrak{J}}\) 是一个双射 \(g\) （从 \(\mathbf{E}^*\) 到 \((\mathrm{E}^{\mathfrak{J}})^*\) 上），它满足 \(g(ax') = a^{\mathfrak{J}}g(x')\) （对一切 \(a \in \mathbf{A}\) ）。因此，\(g\) 与 \((\mathrm{E}^*)^{\mathfrak{J}}\) 到 \(\mathbf{E}^*\) 上的恒同映射的复合，是 \((\mathrm{E}^*)^{\mathfrak{J}}\) 到 \((\mathrm{E}^{\mathfrak{J}})^*\) 上的一个同构。我们将借助这个同构把这些模等同起来，并记作 \(\mathbf{E}_j^*\) 。

设 \(E_{i}, F_{i} (i = 1, \ldots, m)\) 是一些 A-模，\(\Phi_{i} (i = 1, \ldots, m)\) 是 \(E_{i} \times F_{i}\) 上的（关于 J 的）半双线性型。映射

\[
(x _ {1}, \dots , x _ {m}, y _ {1}, \dots , y _ {m}) \rightarrow \Phi_ {1} (x _ {1}, y _ {1}) \Phi_ {2} (x _ {2}, y _ {2}) \dots \Phi_ {m} (x _ {m}, y _ {m})
\]

\((x_{i} \in \mathrm{E}_{i}, y_{i} \in \mathrm{F}_{i}, i = 1, \ldots, m)\) 是 \(\mathrm{E}_{1} \times \cdots \times \mathrm{E}_{m} \times \mathrm{F}_{1}^{J} \times \cdots \times \mathrm{F}_{m}^{J}\) 到 A 中的一个 A-重线性映射，从而定义了 \((\bigotimes_{i} \mathrm{E}_{i}) \times (\bigotimes_{i} \mathrm{F}_{i}^{J})\) 上的一个双线性型（第 III 章，第 1 节，第 7 小节）。由于第二个因子已与 \((\bigotimes_{i} \mathrm{F}_{i})^{J}\) 等同，我们便定义了一个（关于 J 的）半双线性型 \(\Phi\) （在 \((\bigotimes_{i} \mathrm{E}_{i}) \times (\bigotimes_{i} \mathrm{F}_{i})\) 上）。它由下式刻画：

\[
\Phi (x _ {1} \otimes \dots \otimes x _ {m}, y _ {1} \otimes \dots \otimes y _ {m}) = \prod_ {i = 1} ^ {m} \Phi_ {i} (x _ {i}, y _ {i}) \quad (x _ {i} \in \mathrm{E} _ {i}, y _ {i} \in \mathrm{F} _ {i}).\tag{35}
\]

定义 11. --- 给定 A-模 \(\mathbf{E}_{i}, \mathbf{F}_{i}\) （ \(i=1,\ldots,m\) ），并对每个 \(i\) 给定一个（关于 J 的）半双线性型 \(\Phi_{i}\) （在 \(\mathbf{E}_{i} \times \mathbf{F}_{i}\) 上），由 (35) 刻画的半双线性型 \(\Phi\) （关于 J，在 \((\bigotimes_{i}\mathbf{E}_{i}) \times (\bigotimes_{i}\mathbf{F}_{i})\) 上），称为诸半双线性型 \(\Phi_{i}\) 的张量积。

在诸 \(E_{i}\) 与诸 \(F_{i}\) 都等于同一个模 E 且诸 \(\Phi_{i}\) 都等于同一个型 \(\Psi\) 的情形，称 \(\Phi\) 为 \(\Psi\) 到 \(\otimes^{m}E\) 上的扩张。

在定义 11 的记号下，我们来研究与 \(\Phi\) 相伴的映射。由公式 (24)（第 6 小节）与 (35)，我们导出关系式

\[
\Phi (x _ {1} \otimes \dots \otimes x _ {m}, y _ {1} \otimes \dots \otimes y _ {m}) = \prod_ {i = 1} ^ {m} \left\langle x _ {i}, d _ {\Phi_ {i}} (y _ {i}) \right\rangle = \prod_ {i = 1} ^ {m} \left\langle y _ {i}, s _ {\Phi_ {i}} (x _ {i}) \right\rangle^ {j}.
\]

因此我们有：

\[
s _ {\Phi} = j _ {s} \circ (s _ {\Phi_ {1}} \otimes \dots \otimes s _ {\Phi_ {m}}), \quad d _ {\Phi} = j _ {d} \circ (d _ {\Phi_ {1}} \otimes \dots \otimes d _ {\Phi_ {m}})\tag{36}
\]

其中 \(j_s\) （相应地 \(j_d\) ）表示 \(\bigotimes_i F_i^*\) 到 \((\bigotimes_i F_i)^*\) 中（相应地 \(\bigotimes_i E_i^*\) 到 \((\bigotimes_i E_i)^*\) 中）的典范映射（第 III 章，第 1 节，第 4 与第 7 小节）。

命题 9. — 设 A 是一个配备了自同构 J 的交换域，E \(_{i}\) 、F \(_{i}\) 是 A 上两个有限维向量空间，Φ \(_{i}\) 是 E \(_{i}\) × F \(_{i}\) 上的一个（关于 J 的）半双线性型（1 ≤ i ≤ m）。如果诸型 Φ \(_{i}\) 都非退化，则它们的张量积 Φ 也非退化。这时 \(\hat{\Phi}\) （\(\Phi\) 的逆型）是诸逆型 Φ \(_{i}\) 的张量积。

事实上，由于 A 是一个体，由第 III 章，第 1 节，第 3 小节的命题 6 与命题 7 可知，A-模的单射（相应地满射）线性映射的张量积是单射（相应地满射）线性映射。由于诸 \(s_{\Phi_i}\) 依假设是双射（命题 6，第 6 小节），故它们的张量积也是。另一方面，典范映射 \(j_s\) （从 \(\bigotimes_i F_i^*\) 到 \((\bigotimes_i F_i)^*\) 中）是双射（第 III 章，第 1 节，第 5 小节，命题 11）。因此，鉴于 (36)，\(s_\Phi\) 是双射，这就建立了我们的第一个断言（命题 6，第 6 小节）。同样，\(d_\Phi\) 也是双射。

在第二个断言中，我们已把 \(\bigotimes_{i} \mathrm{F}_{i}^{*}\) 与 \((\bigotimes_{i} \mathrm{F}_{i})^{*}\) 、\(\bigotimes_{i} \mathrm{E}_{i}^{*}\) 与 \((\bigotimes_{i} \mathrm{E}_{i})^{*}\) 借助映射 \(j_{s}\) 与 \(j_{d}\) （它们在此是同构）隐式地等同起来。陈述中所引用的诸逆型存在，因为诸 \(s_{\Phi_{i}}\) 、诸 \(d_{\Phi_{i}}\) 、\(s_{\Phi}\) 与 \(d_{\Phi}\) 都是双射（第 7 小节）。令 \(x' = x_{1}' \otimes \cdots \otimes x_{m}'\) ，\(y' = y_{1}' \otimes \cdots \otimes y_{m}'\) （ \(x_{i}' \in \mathrm{E}_{i}^{*}\) ，\(y_{i}' \in \mathrm{F}_{i}^{*}\) ，\(i = 1, \ldots, m\) ）。由逆型的定义，并鉴于 (36)，我们有

\[
\begin{array}{l} \widehat {\Phi} (j _ {s} (y ^ {\prime}), j _ {d} (x ^ {\prime})) = \Phi (s _ {\Phi_ {1}} ^ {- 1} (y _ {1} ^ {\prime}) \otimes \dots \otimes s _ {\Phi_ {m}} ^ {- 1} (y _ {m} ^ {\prime}), d _ {\Phi_ {1}} ^ {- 1} (x _ {1} ^ {\prime}) \otimes \dots \otimes d _ {\Phi_ {m}} ^ {- 1} (x _ {m} ^ {\prime})) \\ = \prod_ {i = 1} ^ {m} \Phi_ {i} (s _ {\Phi_ {i}} ^ {- 1} (y _ {i} ^ {\prime}), d _ {\Phi_ {i}} ^ {- 1} (x _ {i} ^ {\prime})) = \prod_ {i = 1} ^ {m} \widehat {\Phi} _ {i} (y _ {i} ^ {\prime}, x _ {i} ^ {\prime}), \end{array}
\]

由此即得我们的第二个断言。

证毕。

设 E 与 F 是交换环 A 上两个模，\(\Phi\) 是 E × F 上的一个（关于 J 的）半双线性型。映射

\[
(x _ {1}, \dots , x _ {m}, y _ {1}, \dots , y _ {m}) \rightarrow \det (\Phi (x _ {i}, y _ {k})) \quad (x _ {i} \in \mathrm{E}, y _ {i} \in \mathrm{F}, i = 1, \dots , m)
\]

是 \(\mathbf{E}^{m}\times(\mathbf{F}^{\mathrm{j}})^{m}\) 到 A 中的 A-重线性映射。因此它定义了一个双线性型 \(\Phi'\) （在 \((\bigotimes^{\overline{m}}\mathbf{E})\times(\bigotimes^{\overline{m}}\mathbf{F}^{\mathrm{j}})\) 上），由下式刻画

\[
\Phi^ {\prime} (x _ {1} \otimes \dots \otimes x _ {m}, y _ {1} \otimes \dots \otimes y _ {m}) = \det (\Phi (x _ {i}, y _ {k})).
\]

由于当 \(x_{i}=x_{k}\) 或 \(y_{i}=y_{k}\) （ \(i\neq k\) ）时第一项为零，\(\Phi^{\prime}\) 通过过渡到商，定义了 \((\bigwedge^{m}\mathrm{E})\times(\bigwedge^{m}\mathrm{F}^{\mathrm{J}})\) 上的一个双线性型，或者再说一遍，由于 \(\bigwedge^{m}\mathrm{F}^{\mathrm{J}}\) 与 \((\bigwedge^{m}\mathrm{F})^{\mathrm{J}}\) 等同，它定义了一个（关于 J 的）半双线性型 \(\Phi_{(m)}\) （在 \((\bigwedge^{m}\mathrm{E})\times(\bigwedge^{m}\mathrm{F})\) 上）。它由下式刻画

\[
\left\{ \begin{array}{l} \Phi_ {(m)} (x _ {1} \wedge \dots \wedge x _ {m}, y _ {1} \wedge \dots \wedge y _ {m}) = \det (\Phi (x _ {i}, y _ {k})) \\ (x _ {i} \in \mathrm{E}, y _ {i} \in \mathrm{F}, i = 1, \dots , m). \end{array} \right.\tag{37}
\]

定义 12. --- 给定两个 A-模 E、F 与一个半双线性型 \(\Phi\) （关于 J，在 \(E \times F\) 上），由 (37) 刻画的半双线性型 \(\Phi_{(m)}\) （关于 J，在 \((\bigwedge^{m} E) \times (\bigwedge^{m} F)\) 上）称为 \(\Phi\) 到 m 次外幂上的扩张。

在定义 12 的记号下，我们来研究与 \(\Phi_{(m)}\) 相伴的映射。由公式 (24)（第 6 小节）与 (37)，我们导出关系式

\[
\begin{array}{r l} \Phi_ {(m)} (x _ {1} \wedge \dots \wedge x _ {m}, y _ {1} \wedge \dots \wedge y _ {m}) & = \det (\langle x _ {i}, d _ {\Phi} (y _ {k}) \rangle) \\ & = \det (\langle y _ {i}, s _ {\Phi} (x _ {k}) \rangle^ {j}). \end{array}
\]

因此我们有

\[
s _ {\Phi (m)} = k _ {s} \circ (\stackrel {m} {\wedge} s _ {\Phi}), \quad d _ {\Phi_ {m}} = k _ {d} \circ (\stackrel {m} {\wedge} d _ {\Phi}),\tag{38}
\]

其中 \(k_s\) （相应地 \(k_d\) ）表示 \(\bigwedge^m F^*\) 到 \((\bigwedge^m F)^*\) 中（相应地 \(\bigwedge^m E^*\) 到 \((\bigwedge^m E)^*\) 中）的典范映射（参看第 III 章，第 8 节，第 2 小节）。

命题 10. --- 设 A 是一个配备了自同构 J 的交换域，E 与 F 是 A 上两个有限维向量空间，\(\Phi\) 是 E \(\times\) F 上的一个（关于 J 的）半双线性型。若 \(\Phi\) 非退化，则它到 m 次外幂上的扩张 \(\Phi_{(m)}\) 非退化，且 \(\Phi_{(m)}\) 的逆型是逆型 \(\hat{\Phi}\) （\(\Phi\) 的逆型）到 m 次外幂上的扩张。

事实上，由于 \(s_{\Phi}\) 与 \(d_{\Phi}\) 按假设是双射（命题 6，第 6 小节），它们的外幂也是双射（第 III 章，第 5 节，第 7 小节）。另一方面，典范映射 \(k_s\) 与 \(k_d\) 是双射（第 III 章，第 8 节，第 2 小节，定理 1）。于是，鉴于 (38)，\(s_{\Phi(m)}\) 与 \(d_{\Phi(m)}\) 是双射，这就证明了 \(\Phi_{(m)}\) 非退化（命题 6，第 6 小节）。在第二个断言中，我们已把 \(\bigwedge^m F^*\) 与 \((\bigwedge^m F)^*\) 、\(\bigwedge^m E^*\) 与 \((\bigwedge^m E)^*\) 借助映射 \(k_s\) 与 \(k_d\) （它们在此是同构（同前））隐含地等同起来。由于 \(s_{\Phi}\) 、\(d_{\Phi}\) 、\(s_{\Phi(m)}\) 与 \(d_{\Phi(m)}\) 都是双射（第 7 小节），命题中所述的逆型都存在。于是令 \(x' = x_1' \wedge \ldots \wedge x_m'\) 与 \(y' = y_1' \wedge \ldots \wedge y_m'\) （ \(x_i' \in E^*\) ， \(y_i' \in F^*\) ， \(i = 1, \ldots, m\) ）。由逆型的定义（第 7 小节）并鉴于 (38)，我们有

\[
\begin{array}{r l} \hat {\Phi} _ {(m)} (k _ {s} (y ^ {\prime}), k _ {d} (x ^ {\prime})) & = \Phi_ {(m)} (s _ {\Phi} ^ {- 1} (y _ {1} ^ {\prime}) \wedge \dots \wedge s _ {\Phi} ^ {- 1} (y _ {m} ^ {\prime}), d _ {\Phi} ^ {- 1} (x _ {1} ^ {\prime}) \wedge \dots \wedge d _ {\Phi} ^ {- 1} (x _ {m} ^ {\prime})) \\ & = \det (\Phi (s _ {\Phi} ^ {- 1} (y _ {i} ^ {\prime}), d _ {\Phi} ^ {- 1} (x _ {k} ^ {\prime}))) = \det (\hat {\Phi} (y _ {i} ^ {\prime}, x _ {k} ^ {\prime})) \end{array}
\]

由此得到我们的第二个断言。

注记. --- 设 E 是一个自由 A-模，θ 是 ∧E 到 m 阶反称张量所成的子模上的典范同构（第 III 章，第 5 节，第 6 小节，命题 6）。设 Φ 是 E 上的一个半双线性型，Φ(m) 是 Φ 到 ∧E 上的扩张，而 Θ 是 ∧E 上的这样一个半双线性型：它是 Φ 到 ⊗E 上的扩张关于 θ 的原像。由 θ 与张量的反称化的定义以及 (35)，我们有

\[
\Theta (x _ {1} \wedge \dots \wedge x _ {m}, y _ {1} \wedge \dots \wedge y _ {m}) = \sum_ {\sigma , \tau} \varepsilon_ {\sigma} \varepsilon_ {\tau} \Phi (x _ {\sigma (1)}, y _ {\tau (1)}) \dots \Phi (x _ {\sigma (m)}, y _ {\tau (m)})
\]

其中 \(\sigma\) 与 \(\tau\) 取遍对称群 \(\mathfrak{S}_{m}\) 。由行列式的计算公式与公式 (37)，这个表达式可写成

\[
\sum_ {\tau \in \mathfrak {S} _ {m}} \varepsilon_ {\tau} \det (\Phi (x _ {i}, y _ {\tau (k)})) = m! \det (\Phi (x _ {i}, y _ {k}));
\]

换句话说，我们有 \(\Theta = m!\Phi_{(m)}\) 。

